\begin{table}[!htbp]
\centering
\caption{Largest leave-one-family-out movements for the conditioned Q5 target}\label{tab:app_lofo}
\resizebox{\textwidth}{!}{%
\begin{tabular}{rlrrr}
\toprule
Omitted SOC2 & Family & Preperiod stock (\%) & Re-estimated coefficient & Movement [95\% paired CI] \\
\midrule
43 & Office and Administrative Support & 11.21 & -0.0614 & -0.0397 [-0.1131,0.0337] \\
35 & Food Preparation and Serving Related & 3.27 & 0.0174 & 0.0391 [-0.0292,0.1074] \\
31 & Healthcare Support & 3.66 & -0.0534 & -0.0317 [-0.0772,0.0139] \\
13 & Business and Financial Operations & 6.22 & 0.0099 & 0.0316 [-0.0006,0.0638] \\
39 & Personal Care and Service & 2.51 & -0.0457 & -0.0240 [-0.0496,0.0015] \\
25 & Educational Instruction and Library & 4.53 & -0.0351 & -0.0134 [-0.0383,0.0115] \\
17 & Architecture and Engineering & 2.28 & -0.0339 & -0.0123 [-0.0330,0.0085] \\
33 & Protective Service & 2.28 & -0.0096 & 0.0121 [-0.0050,0.0291] \\
\bottomrule
\end{tabular}}
\tabnote{The parent SOC2-by-calendar-month coefficient is $-0.0217$.  Families are ranked by the absolute coefficient movement, so this is an outcome-informed influence diagnostic rather than a preferred deletion rule.  Quintiles and the continuous centering scale remain fixed; all 22 deletions are available in the machine-readable file.}
\end{table}
